\documentclass[UTF8,fontset=none,11pt,a4paper]{ctexart}
\usepackage{ipara-handbook}
\handbooksetup{NET-06 传输层与 TCP}{408 · 计算机网络 NET-06}{Endpoint · UDP · TCP · State}
\providecommand{\tightlist}{\setlength{\itemsep}{0pt}\setlength{\parskip}{0pt}}
\providecommand{\passthrough}[1]{#1}
\begin{document}
\handbooktitle{传输层、端点与 TCP}{从主机交付到进程间有状态通信}{408 · NET-06}
\section{0. 本册定位}\label{0-ux672cux518cux5b9aux4f4d}

本 Topic 回答：IP 只能把 packet 尽力送到 host，怎样进一步把数据交给正确
process，并在需要时维持一条双向、有序、流量受控的 byte-stream
connection？

本册拥有 port、socket endpoint、multiplexing/demultiplexing、UDP、TCP
connection identity、byte sequence state、connection lifecycle、flow
control。

通用 Seq/ACK/Timer/Window
正确性由\href{../03_可靠传输_序号_ACK_定时器与滑动窗口/README.md}{可靠传输}拥有；TCP
拥塞算法由\href{../07_拥塞_共享资源与反馈控制/README.md}{拥塞控制}拥有。本册只解释它们怎样接入
TCP 状态机。

\section{1. 根本问题：IP
地址只能找到机器}\label{1-ux6839ux672cux95eeux9898ip-ux5730ux5740ux53eaux80fdux627eux5230ux673aux5668}

一台 host 同时运行浏览器、DNS resolver、邮件客户端和多个 server
process。只知道 destination IP 仍无法决定把 payload
交给谁。于是传输层增加 endpoint namespace：

\[
\boxed{
\text{Host-to-host delivery}
\xrightarrow{port}
\text{Process endpoint delivery}
}
\]

发送侧 multiplexing 把多个应用的数据交给 UDP/TCP；接收侧 demultiplexing
根据 protocol 与 endpoint identifiers 找到相应 socket/connection state。

\section{2. 名字与对象}\label{2-ux540dux5b57ux4e0eux5bf9ux8c61}

\begin{itemize}
\tightlist
\item
  \textbf{port number}：主机内传输端点的 16-bit 标识空间；
\item
  \textbf{socket endpoint}：至少绑定 local IP、transport protocol、local
  port 等本地信息的通信端点；
\item
  \textbf{TCP connection}：由一对 endpoint 唯一确定，常表示为
  \passthrough{\lstinline!(src IP, src port, dst IP, dst port)!}；
\item
  \textbf{process}：使用 socket API 的执行实体，不等于 port 本身；
\item
  \textbf{segment/datagram}：TCP/UDP 交给 IP 的传输层数据单元。
\end{itemize}

同一 server port 可以同时服务很多 TCP connections，因为 remote endpoint
不同。port 是命名维度，不是``一条连接只能占一个服务器端口''的独占资源。

\section{3.
UDP：保留消息边界的最小传输接口}\label{3-udpux4fddux7559ux6d88ux606fux8fb9ux754cux7684ux6700ux5c0fux4f20ux8f93ux63a5ux53e3}

UDP 增加 source port、destination port、length 和 checksum，然后把一个
application datagram 交给 IP。它不建立 transport
connection，也不维护重传、排序、流控和拥塞窗口状态。

\subsection{3.1 为什么需要
UDP}\label{31-ux4e3aux4ec0ux4e48ux9700ux8981-udp}

朴素方案是所有应用都使用可靠 byte stream。失败在于：

\begin{itemize}
\tightlist
\item
  有些应用需要保留 message boundary；
\item
  有些请求很短，不希望先建立传输连接；
\item
  有些应用要自己控制超时、重传、顺序或实时丢弃策略；
\item
  DHCP 等 bootstrap 流程需要 datagram/broadcast 交互。
\end{itemize}

UDP 不是``更快的
TCP''，而是更少策略的接口。低开销换来的代价是应用必须接受丢失、重复、乱序，或自行构造所需语义。

\subsection{3.2 UDP 首部：四个 16-bit 字段}

UDP 首部固定 8 B：
\[
\boxed{\text{Source Port}\mid\text{Destination Port}\mid\text{Length}\mid\text{Checksum}}.
\]
Length 统计 UDP header 与 data 的总字节数，故最小值为 8；已知 IP total length 与 IP header length 时，也可反推 UDP length。Source Port 在 IPv4 经典规范中可为 0 表示未使用，但 Destination Port 决定接收端服务入口。

Checksum 不只覆盖 UDP header 和 data，还概念性覆盖来自 IP 的 pseudo-header：source/destination IP、上层 protocol/next-header 与 UDP length。伪首部不实际随 UDP 重复传输，它把“这份数据应属于哪对网络端点、哪个上层协议”纳入完整性检查，降低误投递未被发现的风险。IPv4 UDP checksum 可选与 IPv6 中通常必需的边界必须服从题设，不能把 IPv4 的全零语义无条件搬到 IPv6。

\section{TCP 首部：连接状态在报文中的投影}

TCP 基础首部最少 20 B，可压缩为：
\[
\boxed{\begin{gathered}
\text{Src/Dst Port}\mid\text{Sequence Number}\mid\text{Acknowledgment Number}\\
\text{Data Offset/Flags}\mid\text{Window}\mid\text{Checksum}\mid\text{Urgent Pointer}\\
\text{Options+Padding}\mid\text{Data}
\end{gathered}}
\]

字段应按“它暴露哪份状态”来读：
\begin{itemize}
\tightlist
\item port pair 与 IP pair 一起选择 connection；
\item SEQ 标记本段首个数据字节，ACK 在 ACK flag 有效时表示下一期望字节；
\item Data Offset 以 32-bit word 为单位给出 TCP header length，最小为 5，故最小首部 $5\times4=20$ B；
\item Window 表示从 ACK number 开始、该发送者当前愿意接收的字节数，是接收方向对反向发送者发布的流控状态；
\item Checksum 覆盖 pseudo-header、TCP header 与 data，TCP 中不是可选项；
\item Urgent Pointer 只有 URG 置位才有意义；Options 由 Data Offset 划定，MSS 是建连时常见选项而非固定基础字段。
\end{itemize}

408 常用 flags 可用状态事件复原：SYN 同步序号，ACK 使 acknowledgment 字段有效，FIN 表示本方向不再发送，RST 异常复位/拒绝，PSH 请求及时向应用推进，URG 使紧急指针有效。SYN 与 FIN 各占一个序号；纯 ACK 若不带数据且无 SYN/FIN，通常不消耗序号。字段存在不表示语义总生效，例如 ACK number 只有 ACK flag 置位才应解释。

TCP/UDP checksum 的 pseudo-header 是跨层校验输入，不是“传输层首部的一部分”，也不由中间路由器作为 TCP/UDP 字段逐跳改写。NAT 改写 IP/port 后必须相应维护校验和，正是因为这些值参与端点完整性约束。

\section{5. TCP：把不可靠 packet 网络解释成可靠 byte
stream}\label{4-tcpux628aux4e0dux53efux9760-packet-ux7f51ux7edcux89e3ux91caux6210ux53efux9760-byte-stream}

TCP 向应用提供有序 byte stream，而不是保留应用 write
的消息边界。每个传输方向都有独立序号空间和状态：

\[
\boxed{
\text{Application bytes}
\to \text{Byte sequence space}
\to \text{Segments}
\to \text{ACKed contiguous prefix}
\to \text{Receive byte stream}
}
\]

TCP sequence number 标记 segment 中第一个 data octet 的序号；ACK
通常表示``下一期望 byte''，也就是此前连续 byte 已收到。SYN 和 FIN 也占用
sequence space，这使控制事件能与数据一起被可靠排序。

\section{6.
为什么连接建立需要同步双方状态}\label{5-ux4e3aux4ec0ux4e48ux8fdeux63a5ux5efaux7acbux9700ux8981ux540cux6b65ux53ccux65b9ux72b6ux6001}

建立 TCP connection 不是``测试网络通不通''，而是让双方确认：

\begin{itemize}
\tightlist
\item
  对端 endpoint 存在且愿意通信；
\item
  双方 initial sequence number 已知；
\item
  旧连接的迟到 segment 不应被误认作当前 connection 数据；
\item
  双向发送/接收状态可以进入 ESTABLISHED。
\end{itemize}

典型三次握手：

\begin{lstlisting}
Client                         Server
CLOSED                         LISTEN
SYN, seq=x ------------------>
             <--------------- SYN+ACK, seq=y, ack=x+1
ACK, ack=y+1 ---------------->
ESTABLISHED                    ESTABLISHED
\end{lstlisting}

第三次消息的关键不是``为了凑三次''，而是让 server 知道 client 已收到
server 的 SYN/ISN。两次消息只能让 client 知道双向可达，server
仍无法确认其 SYN 已被 client 接受。

\subsection{5.1
状态比报文名更重要}\label{51-ux72b6ux6001ux6bd4ux62a5ux6587ux540dux66f4ux91cdux8981}

主动打开常见轨迹：\(\text{CLOSED} \rightarrow \text{SYN-SENT} \rightarrow \text{ESTABLISHED}\)。

被动打开常见轨迹：\(\text{CLOSED} \rightarrow \text{LISTEN} \rightarrow \text{SYN-RECEIVED} \rightarrow \text{ESTABLISHED}\)。

同时打开、半开连接、RST 和重传会产生其他分支。做题必须用``收到什么
segment 前处于什么 state''推演，而不是只背三行握手图。

\section{7. TCP
数据传输的三套约束}\label{6-tcp-ux6570ux636eux4f20ux8f93ux7684ux4e09ux5957ux7ea6ux675f}

发送方在任意时刻可发送多少数据，至少同时受：

\begin{enumerate}
\def\labelenumi{\arabic{enumi}.}
\tightlist
\item
  \textbf{可靠性状态}：哪些 bytes 已发送、已确认、允许重传；
\item
  \textbf{flow control}：receiver advertised window
  \passthrough{\lstinline!rwnd!}；
\item
  \textbf{congestion control}：sender 的
  \passthrough{\lstinline!cwnd!}。
\end{enumerate}

教材常压缩为：

\[
W_{usable}\le \min(rwnd,cwnd)-\text{bytes in flight}.
\]

\passthrough{\lstinline!rwnd!} 来自接收端 buffer/应用读取状态，保护
receiver；\passthrough{\lstinline!cwnd!} 来自 sender 对 network path
的拥塞推断，保护
network。二者数值都像``窗口''，但状态所有者、反馈来源和失败对象不同。

\section{8. Flow
control：接收端怎样表达``我还能接多少''}\label{7-flow-controlux63a5ux6536ux7aefux600eux6837ux8868ux8fbeux6211ux8fd8ux80fdux63a5ux591aux5c11}

接收端将可用 buffer 空间通过 advertised window
反馈给发送端。应用读取数据会释放空间；新数据到达会占用空间。基本不变量是：发送方不能让未消费数据超过接收方声明的可接受范围。

若 receiver 通告 zero window，sender
暂停正常新数据发送，但必须保留恢复探测机制，否则 window update
丢失可能造成双方永久等待。

Flow control 不证明数据已可靠交付，也不代表路径没有拥塞。它只约束
receiver capacity。

\section{9. TCP
可靠性如何实例化通用机制}\label{8-tcp-ux53efux9760ux6027ux5982ux4f55ux5b9eux4f8bux5316ux901aux7528ux673aux5236}

TCP 使用 byte sequence number、ACK、retransmission timer 和 duplicate
suppression，把\href{../03_可靠传输_序号_ACK_定时器与滑动窗口/README.md}{可靠传输}的机制用于
byte stream。

但 TCP 不能简单等同 GBN 或 SR：

\begin{itemize}
\tightlist
\item
  ACK 通常具有 cumulative semantics；
\item
  receiver 可以缓存 out-of-order bytes；
\item
  retransmission 可以由 timeout 或 duplicate ACK/SACK 等信号触发；
\item
  具体实现与扩展比教材中的纯 GBN/SR 状态机更复杂。
\end{itemize}

因此 GBN/SR 是理解窗口正确性和代价的模型，不是给 TCP 贴唯一协议标签。

\section{10.
连接终止：两个方向必须分别关闭}\label{9-ux8fdeux63a5ux7ec8ux6b62ux4e24ux4e2aux65b9ux5411ux5fc5ux987bux5206ux522bux5173ux95ed}

TCP 是 full-duplex。一个 endpoint
不再发送，并不意味着它也不能继续接收。典型主动关闭：

\begin{lstlisting}
A sends FIN
-> B ACKs: A->B direction closed after prior bytes
-> B may continue sending
-> B sends FIN when its direction finishes
-> A ACKs
-> active closer waits in TIME-WAIT before state disappears
\end{lstlisting}

四次报文来自两个方向的关闭意图可能不同时发生；若 ACK 与 FIN
可合并，报文数可能少于四个。不能把``四次挥手''当作固定物理定律。

\subsection{9.1 TIME-WAIT
维护什么}\label{91-time-wait-ux7ef4ux62a4ux4ec0ux4e48}

TIME-WAIT 让 active closer 有机会重发最后 ACK，并让旧 connection 的迟到
segment 在相同 endpoint tuple
被复用前消失。它用暂时保留状态换取连接实例之间的不歧义。

\section{11. TCP
状态的一生}\label{10-tcp-ux72b6ux6001ux7684ux4e00ux751f}

每条 connection 至少维护：

\begin{itemize}
\tightlist
\item
  local/remote endpoint；
\item
  send/receive initial and current sequence variables；
\item
  send/receive window；
\item
  retransmission and timing state；
\item
  connection state；
\item
  congestion-control state（由拥塞 Topic 解释）。
\end{itemize}

\begin{lstlisting}
Create endpoint
-> Open / Listen
-> Synchronize sequence spaces
-> Transfer bytes in both directions
-> Retransmit / reorder / flow-control as events occur
-> Half-close or abort
-> Release state after safety interval
\end{lstlisting}

TCP 的代价不是只多 20-byte header，而是 endpoint 长期持有并更新这些
per-connection states。

\section{12. 408 协议流程：UDP/TCP 逐字段推进}

\subsection{UDP demultiplexing 与 checksum}

发送应用产生一个 datagram；UDP 加入 source port、destination port、length 和 checksum 后交给 IP。接收端先按 IP protocol 找到 UDP，再验证 length/checksum，最后以 local address/port 等绑定状态选择 endpoint；无匹配 endpoint 时不能把 payload 随意交给某个 process。UDP checksum 检查通过只保护报文完整性，不产生重传、排序或连接状态。

\subsection{三次握手：同步两个序号空间}

\begin{handbooktable}{L{2.0cm} L{3.5cm} Y Y}
\toprule
步骤 & segment & Client state / knowledge & Server state / knowledge\\
\midrule
1 & SYN, seq=$x$ & CLOSED $\to$ SYN-SENT；占用 $x$ & LISTEN 收到 client ISN\\
2 & SYN+ACK, seq=$y$, ack=$x+1$ & 知道 server 收到 $x$ 且获知 $y$ & SYN-RECEIVED；等待 $y$ 被确认\\
3 & ACK, ack=$y+1$ & ESTABLISHED & 确认 client 已收到 server ISN；ESTABLISHED\\
\bottomrule
\end{handbooktable}

SYN 重传、同时打开和 RST 是分支；正常三行不是所有输入下的唯一轨迹。每一步都必须依据收到的 flag、SEQ、ACK 和当前 state 判断是否合法。

\subsection{数据、ACK、超时与快速重传接口}

TCP SEQ 指向本 segment 首个 data octet；ACK 指向下一期望 byte。收到乱序 segment 时可缓存并重复发送累计 ACK；连续缺口被填上后 ACK 才跨过该段。RTO 到期触发可靠性重传，并向 NET07 提供强拥塞信号；多个 duplicate ACK 可触发 fast retransmit，并把拥塞窗口分支交给 NET07。这里 Owns TCP 报文与连接状态，NET03 Owns通用正确性，NET07 Owns `cwnd` 更新。

\subsection{接收窗口与 zero-window probe}

接收方通告 `rwnd = receive buffer capacity - buffered unread bytes` 的题设模型。发送方正常可用量为
\[
\max\{0,\min(rwnd,cwnd)-FlightSize\}.
\]
若收到 zero window，发送方暂停正常新数据，但保留 persist/probe 分支；探测促使接收方重报当前窗口，避免窗口更新丢失后双方永久等待。该机制保护 receiver，不证明 network 未拥塞。

\subsection{关闭、半关闭、TIME-WAIT 与 RST}

FIN 只关闭一个发送方向，并占用一个 sequence number；对端 ACK 后仍可在反方向发送。主动关闭方确认对端 FIN 后进入 TIME-WAIT，以便重发最终 ACK 并隔离旧 segment。RST 表示拒绝、异常终止或不存在的连接状态，不执行有序半关闭。题目若把 ACK 与 FIN 合并，报文数可以少于经典四段，但两个方向的关闭状态仍必须分别核对。

\section{13. 概念边界}\label{11-ux6982ux5ff5ux8fb9ux754c}

\begin{longtable}[]{@{}lclll@{}}
\toprule\noalign{}
概念 A & ≠ & 概念 B & 真正区别与题目信号 & 混淆后果 \\
\midrule\noalign{}
\endhead
\bottomrule\noalign{}
\endlastfoot
Port & ≠ & Process & port 是 transport namespace；process 是使用
endpoint 的执行实体 & 认为一个 port 只能对应一次会话 \\
Socket endpoint & ≠ & TCP connection & endpoint 是一端；connection
由两端共同确定 & 四元组与本地绑定混乱 \\
UDP datagram & ≠ & TCP byte stream & UDP 保留 message boundary；TCP
只保证 byte order & 应用分帧假设错误 \\
Connection establishment & ≠ & Data reliability & 握手同步 connection
state；数据可靠性需持续 ACK/重传 & 认为握手后不会丢包 \\
ACK number & ≠ & Segment number & TCP ACK/SEQ 以 byte space 为核心 & 按
packet 个数推进序号 \\
\passthrough{\lstinline!rwnd!} & ≠ & \passthrough{\lstinline!cwnd!} &
receiver capacity 与 network capacity & 流控/拥塞触发和更新写反 \\
FIN & ≠ & RST & FIN 有序关闭一个方向；RST 异常终止/拒绝状态 &
关闭状态机推演错误 \\
TIME-WAIT & ≠ & ``server 固定等待'' & 通常由 active closer
进入，保护最终 ACK 与旧 segment 隔离 & 只按 client/server 身份判断 \\
Pseudo-header & ≠ & Transmitted transport header & 前者只参与 checksum 计算，字段来自 IP & 错算首部长度或线速负载 \\
UDP Length & ≠ & IP Total Length & 前者含 UDP header+data；后者还含 IP header & 反推 payload 时漏减首部 \\
\end{longtable}

\section{14.
做题调用协议}\label{12-ux505aux9898ux8c03ux7528ux534fux8bae}

\begin{enumerate}
\def\labelenumi{\arabic{enumi}.}
\tightlist
\item
  写 transport protocol 和完整 endpoint tuple；
\item
  明确应用需要 datagram 还是 byte stream 语义；
\item
  TCP 题为两个方向分别画
  \passthrough{\lstinline!SEQ/ACK/state/window!}；
\item
  先确定 SYN/FIN 是否占序号，再计算 ACK；
\item
  窗口题分开 bytes ACKed、in
  flight、\passthrough{\lstinline!rwnd!}、\passthrough{\lstinline!cwnd!}；
\item
  建连/关闭题逐事件更新 state，不背固定报文数；
\item
  最后检查同一 server port 的不同 connections 是否被误合并。
\item
  首部题先用 Data Offset/Length 划定边界，再按 flags 判断 SEQ、ACK、Window、Urgent Pointer 是否有效；
\item
  checksum 题把 pseudo-header 列为计算输入，但不得把它计入实际传输首部长度。
\end{enumerate}

\section{15. 贯穿母例：同一 80
端口为什么能服务多个浏览器}\label{13-ux8d2fux7a7fux6bcdux4f8bux540cux4e00-80-ux7aefux53e3ux4e3aux4ec0ux4e48ux80fdux670dux52a1ux591aux4e2aux6d4fux89c8ux5668}

server 在 \passthrough{\lstinline!203.0.113.8:80!} listen。两个 clients
建立：

\begin{lstlisting}
(192.0.2.10:51000, 203.0.113.8:80)
(198.51.100.20:52000, 203.0.113.8:80)
\end{lstlisting}

server local port 相同，但 remote IP/port 不同，因此是两条独立
connection，各自拥有 sequence、window、timer 和 state。若只看
destination port，会把多条状态机错误地合成一条。

\section{16. 高频 First
Divergence}\label{14-ux9ad8ux9891-first-divergence}

\begin{itemize}
\tightlist
\item
  把 port 写成应用程序本身：没有区分 namespace 与 owner；
\item
  TCP sequence 每 segment 加 1：忘记它以 byte 为基本坐标，SYN/FIN
  另有占用；
\item
  三次握手解释为``第三次确认数据''：没有说明 server ISN 的确认闭环；
\item
  看到窗口只取 \passthrough{\lstinline!rwnd!} 或只取
  \passthrough{\lstinline!cwnd!}：漏掉共同约束和 in-flight 数据；
\item
  把 TCP 判成 GBN：把教材可靠协议类比当成协议事实；
\item
  默认 client 总进入 TIME-WAIT：没有看谁主动关闭。
\item
  把 Data Offset 当字节数：忘记其单位是 32-bit word；
\item
  把 pseudo-header 画进在线 TCP/UDP 首部：混淆校验输入与传输字段；
\item
  看到 ACK number 就推进确认：没有先检查 ACK flag 与当前连接状态。
\end{itemize}

\section{17.
一页压缩与复原问题}\label{15-ux4e00ux9875ux538bux7f29ux4e0eux590dux539fux95eeux9898}

\[
\boxed{
\text{IP host delivery}
\to \text{Port demultiplexing}
\to \text{Endpoint pair}
\to \text{Sequence-space synchronization}
\to \text{Bidirectional byte-state evolution}
}
\]

\begin{enumerate}
\def\labelenumi{\arabic{enumi}.}
\tightlist
\item
  为什么 port 不能唯一标识一条 TCP connection？
\item
  UDP 少掉的状态给应用带来什么自由和责任？
\item
  三次握手的第三次究竟让谁知道了什么？
\item
  \passthrough{\lstinline!rwnd!}、\passthrough{\lstinline!cwnd!} 和
  reliable sequence state 分别保护什么？
\item
  TIME-WAIT 怎样避免新旧 connection 混淆？
\item
  TCP/UDP 的 pseudo-header 为什么能发现一部分误投递，又为什么不增加线上首部长度？
\end{enumerate}

\section{18.
来源与校正说明}\label{16-ux6765ux6e90ux4e0eux6821ux6b63ux8bf4ux660e}

\begin{itemize}
\tightlist
\item
  对应归档中的《传输层-TCP》《传输层-传输层功能与UDP》为空文件，因此未把空入口当作已有模型；
\item
  TCP byte sequence、ACK、connection state 与接口边界依据
  \href{https://www.rfc-editor.org/rfc/rfc9293.html}{RFC 9293} 校正；
\item
  本册保留 408 的握手、关闭、流控与端口模型，同时把现代扩展与纯 GBN/SR
  类比限制在明确边界内。
\end{itemize}
\end{document}
